\documentclass[11pt]{article}
\usepackage[utf8]{inputenc}
\usepackage[spanish]{babel}
\usepackage{geometry}
\geometry{a4paper, margin=2.5cm}
\usepackage{graphicx}
\usepackage{tikz}
\usepackage{pgfplots}
\usetikzlibrary{shapes, arrows, positioning, fit}
\usepackage{xcolor}
\usepackage{hyperref}
\usepackage{booktabs}
\usepackage{array}
\usepackage{float}

% Colores corporativos
\definecolor{corporate-blue}{RGB}{0, 82, 155}
\definecolor{corporate-gray}{RGB}{240, 240, 240}
\definecolor{accent-orange}{RGB}{255, 102, 0}

% Estilo profesional
\usepackage{titlesec}
\titleformat{\section}{\color{corporate-blue}\normalfont\Large\bfseries}{}{0pt}{}
\titleformat{\subsection}{\color{corporate-blue}\normalfont\large\bfseries}{}{0pt}{}

\usepackage{fancyhdr}
\pagestyle{fancy}
\fancyhf{}
\fancyhead[L]{\textcolor{corporate-blue}{\footnotesize Análisis Financiero Corporativo}}
\fancyhead[R]{\textcolor{corporate-blue}{\footnotesize \thepage}}
\renewcommand{\headrulewidth}{0.4pt}
\renewcommand{\headrule}{\color{corporate-blue}\hrule}

\begin{document}

% Portada corporativa
\begin{titlepage}
    \centering
    \vspace*{2cm}
    % \includegraphics[width=0.3\textwidth]{logo_empresa.png} % Reemplazar con logo real
    \vspace{1cm}
    
    {\Huge\color{corporate-blue}\textbf{ANÁLISIS FINANCIERO CORPORATIVO}}
    \vspace{1cm}
    
    {\Large\textbf{Informe de Diagnóstico Integral}}
    \vspace{1.5cm}
    
    {\large Generado el: 18 de October de 2025}
    \vspace{0.5cm}
    
    {\large Documentos analizados: 1}
    \vspace{2cm}
    
    \rule{\textwidth}{0.4pt}
    \vspace{0.5cm}
    
    {\large Departamento de Análisis Financiero}
    \vfill
\end{titlepage}

% Tabla de contenidos
\tableofcontents
\newpage


    % Resumen Ejecutivo
    \section{Resumen Ejecutivo}
    \noindent\fbox{\parbox{\textwidth}{
    Análisis consolidado de 1 documentos financieros. Se identificaron 3 puntos fuertes principales y 2 áreas críticas de mejora.
    }}

    % Análisis por Documento
    \section{Análisis por Documento}
    
    \subsection{manual_empresa.pdf}
    \textbf{Resumen:} TeenFinanKids es una aplicación móvil que brinda educación financiera integral a niños y jóvenes en Colombia.

    \textbf{Riesgo:} bajo

    \subsubsection{Puntos Fuertes}
    \begin{itemize}
        \item aplicación móvil
    \item educación financiera integral
    \item participación de padres y tutores
\end{itemize}

    \subsubsection{Áreas de Mejora}
    \begin{itemize}
        \item limitaciones de implementación
    \item riesgos económicos
\end{itemize}

    \subsubsection{Recomendaciones}
    \begin{itemize}
        \item fomentar hábitos financieros saludables desde la infancia
    \item consolidar una experiencia educativa motivadora
\end{itemize}



% Sección de Grafos y Diagramas
\section{Análisis Gráfico y Diagramas}

% Gráfico de Riesgos
\subsection{Distribución de Riesgos}
\begin{figure}[H]
    \centering
    \begin{tikzpicture}
        \begin{axis}[
            ybar,
            bar width=0.8cm,
            width=0.8\textwidth,
            height=6cm,
            symbolic x coords={Alto, Medio, Bajo},
            xtick=data,
            nodes near coords,
            nodes near coords align={vertical},
            ymin=0,
            ymax=3,
            xlabel={Nivel de Riesgo},
            ylabel={Cantidad de Documentos},
            tick label style={font=\small},
            label style={font=\small}
        ]
        \addplot[fill=corporate-blue] coordinates {(Alto, 0) (Medio, 0) (Bajo, 1)};
        \end{axis}
    \end{tikzpicture}
    \caption{Distribución de niveles de riesgo identificados}
\end{figure}

% Diagrama de Flujo de Procesos
\subsection{Diagrama de Procesos Identificados}

\begin{figure}[H]
    \centering
    \begin{tikzpicture}[node distance=2cm, auto]
        \tikzstyle{process} = [rectangle, draw=corporate-blue, fill=corporate-gray, thick, minimum width=3cm, minimum height=1cm, text centered, rounded corners]
        \tikzstyle{decision} = [diamond, draw=accent-orange, fill=orange!20, thick, minimum width=2cm, minimum height=1cm, text centered, aspect=2]
        \tikzstyle{arrow} = [thick,->,>=stealth]
        
        \node (start) [process] {Inicio del Proceso};
        \node (analysis) [process, below of=start] {Análisis Financiero};
        \node (decision) [decision, below of=analysis] {¿Riesgo Alto?};
        \node (action1) [process, right of=decision, xshift=3cm] {Acción Correctiva};
        \node (action2) [process, below of=decision, yshift=-1cm] {Seguimiento Normal};
        \node (end) [process, below of=action2] {Finalización};
        
        \draw [arrow] (start) -- (analysis);
        \draw [arrow] (analysis) -- (decision);
        \draw [arrow] (decision) -- node[near start] {Sí} (action1);
        \draw [arrow] (decision) -- node[near start] {No} (action2);
        \draw [arrow] (action1) |- (end);
        \draw [arrow] (action2) -- (end);
    \end{tikzpicture}
    \caption{Diagrama de flujo del proceso de análisis financiero}
\end{figure}


% Grafo de Relaciones Departamentales  
\subsection{Mapa de Relaciones Organizacionales}

\begin{figure}[H]
    \centering
    \begin{tikzpicture}[
        node distance=3cm,
        department/.style={rectangle, draw=corporate-blue, fill=corporate-gray, thick, minimum width=2.5cm, minimum height=1cm, text centered, rounded corners},
        relation/.style={->, thick, >=stealth}
    ]
    
        \node[department] (finance) {Finanzas};
        \node[department, right of=finance] (operations) {Operaciones};
        \node[department, below of=finance] (sales) {Ventas};
        \node[department, right of=sales] (marketing) {Marketing};
        
        \draw[relation] (sales) -- node[midway, left] {Informes} (finance);
        \draw[relation] (operations) -- node[midway, above] {Costos} (finance);
        \draw[relation] (marketing) -- node[midway, right] {Presupuesto} (finance);
        \draw[relation] (sales) -- node[midway, below] {Coordinación} (operations);
        
    \end{tikzpicture}
    \caption{Mapa de relaciones interdepartamentales identificadas}
\end{figure}


% Mapa Conceptual de Términos
\subsection{Mapa Conceptual Financiero}

\begin{figure}[H]
    \centering
    \begin{tikzpicture}[
        concept/.style={ellipse, draw=accent-orange, fill=orange!20, thick, minimum width=2cm, minimum height=1cm, text centered},
        connection/.style={->, thick, >=stealth, bend left=15}
    ]
    
        \node[concept] (income) {Ingresos};
        \node[concept, right of=income, xshift=3cm] (costs) {Costos};
        \node[concept, below of=income, yshift=-1cm] (profit) {Utilidades};
        \node[concept, right of=profit, xshift=3cm] (cashflow) {Flujo de Caja};
        
        \draw[connection] (income) to node[midway, above] {Genera} (profit);
        \draw[connection] (costs) to node[midway, above] {Afecta} (profit);
        \draw[connection] (profit) to node[midway, right] {Impacta} (cashflow);
        \draw[connection] (income) to node[midway, left] {Financia} (cashflow);
        
    \end{tikzpicture}
    \caption{Mapa conceptual de términos financieros clave}
\end{figure}


% Arquitectura Organizacional
\subsection{Diagrama de Arquitectura Organizacional}

\begin{figure}[H]
    \centering
    \begin{tikzpicture}[
        level 1/.style={sibling distance=4cm},
        level 2/.style={sibling distance=2cm},
        every node/.style={rectangle, draw=corporate-blue, fill=corporate-gray, thick, minimum width=2.5cm, text centered, rounded corners}
    ]
    
    \node {Gerencia General}
        child {node {Finanzas}
            child {node {Contabilidad}}
            child {node {Análisis}}
        }
        child {node {Operaciones}
            child {node {Producción}}
            child {node {Logística}}
        }
        child {node {Comercial}
            child {node {Ventas}}
            child {node {Marketing}}
        };
    
    \end{tikzpicture}
    \caption{Arquitectura organizacional sugerida}
\end{figure}



\end{document}